\section*{Chapter \ref{ch:g11:public-goods-and-commons} --- Public Goods and Common Resources}
% figdata: figdata/grade-11/07-public-goods-and-commons/out/ (peru_key, peru, itq_key); Kestria numbers invented.

\begin{solution}{exo:g11:public-goods-and-commons:1}
A sandwich: private (rival, excludable). Fish in the open ocean: \omterm{def:g11:public-goods-and-commons:private-good}{common resource} (rival, not excludable). National defence: \omterm{def:g11:public-goods-and-commons:private-good}{public good} (protecting one person does not
protect another less, and no one in the country can be left out). A film in a cinema: \omterm{def:g11:public-goods-and-commons:private-good}{club good} (a ticket is needed, and one more viewer leaves the film to the others,
until the seats run out).
\end{solution}

\begin{solution}{exo:g11:public-goods-and-commons:2}
Fireworks seen from all over town: \omterm{def:g11:public-goods-and-commons:private-good}{public good}. A swimming pool that is rarely full: \omterm{def:g11:public-goods-and-commons:private-good}{club good}. Groundwater: \omterm{def:g11:public-goods-and-commons:private-good}{common resource}, since what one farmer pumps is lost to the
others and no one can be stopped from drilling. Shoes: \omterm{def:g11:public-goods-and-commons:private-good}{private good}.
\end{solution}

\begin{solution}{exo:g11:public-goods-and-commons:3}
Empty at night: a \omterm{def:g11:public-goods-and-commons:private-good}{public good}, since one more car takes nothing from the others. Jammed at rush hour: a \omterm{def:g11:public-goods-and-commons:private-good}{common resource}, since each car slows all the others. A toll makes
the road excludable: a \omterm{def:g11:public-goods-and-commons:private-good}{club good} at night, and at rush hour a \omterm{def:g11:public-goods-and-commons:private-good}{private good}, in effect, whose price can keep the road from jamming.
\end{solution}

\begin{solution}{exo:g11:public-goods-and-commons:4}
12.5 million tonnes in 1970 and 2.3 million in 1973: a fall of $(12.5 - 2.3)/12.5 \approx 82\,\%$.
\end{solution}

\begin{solution}{exo:g11:public-goods-and-commons:5}
In 1994, 21 years after 1973.
\end{solution}

\begin{solution}{exo:g11:public-goods-and-commons:6}
A boat's catch is worth $2(50 - n/2) = 100 - n$; boats enter until $100 - n = 60$, so $n = 40$. Each catches 30 tonnes; the total is 1200.
\end{solution}

\begin{solution}{exo:g11:public-goods-and-commons:7}
$100 - n = 40$ gives $n = 60$; each boat catches 20 tonnes, and the total is again 1200. No one gains: the saving on engines is used up by the 20 extra boats, and the
rent stays zero.
\end{solution}

\begin{solution}{exo:g11:public-goods-and-commons:8}
$R(n) = n(100 - n) - 60n = 40n - n^2$, largest at $n = 20$. Each boat catches 40 tonnes; the total is 800 and the rent $20 \times (80 - 60) = 400$ dalens.
\end{solution}

\begin{solution}{exo:g11:public-goods-and-commons:9}
800 tonnes. At 20 boats a tonne costs $60/40 = 1.5$ dalens to catch and sells for 2, so the right to catch it leases for 0.5 dalen ($400/800$).
\end{solution}

\begin{solution}{exo:g11:public-goods-and-commons:10}
$30\,000/n + 20 = 80$ gives $n = 500$ members. The club can shut its gate to those who do not pay, so a fee can be charged; a \omterm{def:g11:public-goods-and-commons:private-good}{public good} cannot exclude, so a fee cannot
be enforced.
\end{solution}

\begin{solution}{exo:g11:public-goods-and-commons:11}
New Zealand: about 210 thousand tonnes in 1986 and 365 thousand in 2020. Iceland: about 1520 thousand tonnes in 1990 and 1040 thousand in 2020. They do not show it either
way: a catch is set by the quota itself, by the stock and by the market, and a lower catch may mean a stock protected by a tight quota or a stock in decline. The test needs
measures of the stocks, across many fisheries with and without quotas.
\end{solution}

\begin{solution}{exo:g11:public-goods-and-commons:12}
\omterm{def:g11:public-goods-and-commons:private-good}{Common resources} kept in good order for a long time by their users, without private ownership of the resource or state control: the mountain pastures, irrigation systems and
fisheries of the field studies. They break the assumption of \omterm{def:g11:perfect-competition:free-entry}{free entry} (outsiders are kept out, so the resource is not open access) and the assumption that each user ignores
the others' loss (users who watch each other and punish overuse make that loss their own).
\end{solution}

\begin{solution}{exo:g11:public-goods-and-commons:13}
Player $i$ gets $10 - g_i + \frac{m}{n}(g_1 + \dots + g_n)$. A dalen more from $i$ changes this by $-1 + m/n < 0$, whatever the others give, so $g_i = 0$ is dominant. If all
give 10, each gets $\frac{m}{n} \times 10n = 10m > 10$, what each gets if all give nothing.
\end{solution}

\begin{solution}{exo:g11:public-goods-and-commons:14}
\omterm{def:g11:oligopoly-and-games:game}{Payoffs} (first country, second): both cut $(-10 + 12,\ -10 + 12) = (2, 2)$; only the first cuts $(-4, 6)$; only the second cuts $(6, -4)$; neither $(0, 0)$. Not cutting is
dominant for each ($0 > -4$ and $6 > 2$), so the equilibrium is $(0, 0)$, worse for both than $(2, 2)$: a \omterm{def:g11:oligopoly-and-games:prisoners-dilemma}{prisoner's dilemma}. An agreement must be checked and must punish
the country that breaks it, for instance by the others' withdrawing their own cuts in later rounds.
\end{solution}

\begin{solution}{exo:g11:public-goods-and-commons:15}
$6/0.05 = 120$ dalens. With a rent of 8, $8/0.05 = 160$. If the quota ends after one season, it is worth only that season's rent, about 6 dalens.
\end{solution}

\begin{solution}{pb:g11:public-goods-and-commons:1}
\textbf{1.} One boat: $6(100 - n) = 600 - 6n$ kests. The bay: $n(600 - 6n) = 600n - 6n^2$.

\textbf{2.} $600 - 6n = 120$: $n = 80$.

\textbf{3.} 20 tonnes a boat, 1600 tonnes in all, rent $9600 - 9600 = 0$.

\textbf{4.} With 79 boats each catches 21 tonnes; the 80th lowers that to 20, so the other 79 lose 79 tonnes in all, while it catches 20: the total catch falls by 59.

\textbf{5.} The owner pays his boat's cost and receives his boat's catch; the tonne each of the others loses is their loss, not his.

\textbf{6.} $R(n) = 600n - 6n^2 - 120n = 480n - 6n^2$; $R'(n) = 480 - 12n = 0$ at $n = 40$.

\textbf{7.} 60 tonnes a boat, 2400 tonnes, rent $14\,400 - 4800 = 9600$ kests.

\textbf{8.} $n(100 - n)$ is largest at $n = 50$: 2500 tonnes.

\textbf{9.} 50 boats catch most (2500 tonnes); 40 boats earn most (9600 kests); open access, 80 boats, catches least (1600) and earns nothing.

\textbf{10.} 9600 kests a season.

\textbf{11.} 2400 quotas.

\textbf{12.} A boat costs 120 and catches 60 tonnes: 2 kests a tonne. The rent is $6 - 2 = 4$ kests a tonne ($9600/2400$).

\textbf{13.} A boat earns 4 kests above its cost on each tonne it is allowed to catch. A holder will not lease for less, and a fisher will not pay more; with many of each,
the lease settles at 4.

\textbf{14.} $4/0.05 = 80$ kests.

\textbf{15.} The quota's price is the value of the rents it will earn; a healthier stock means cheaper fishing and a higher rent, so a run-down stock lowers the value of
what the holder owns.

\textbf{16.} 30 tonnes each, worth $30 \times 80 = 2400$ kests.

\textbf{17.} Each sells its 30 tonnes of quota, 2400 kests, to the boats that stay, which end with 60 tonnes each.

\textbf{18.} $n(100 - n) = 2100$ gives $n = 30$ or $70$; the cheapest way to catch 2100 tonnes is with 30 boats, each catching 70. The rent is $6 \times 2100 - 120 \times 30 = 9000$,
600 less than at the efficient total.

\textbf{19.} For example: counting each boat's landings and punishing cheating; fish thrown back to keep within the quota; quotas gathered in a few hands.

\textbf{20.} Open access dissipates the whole rent of the bay, 9600 kests a season; under quotas, the right to catch one tonne for one season leases for 4 kests, and a
permanent one-tonne quota sells for 80 kests at 5\,\% interest.
\end{solution}
